\chapter{Encopresis}

\section*{Definition}
Repeated passage of feces into inappropriate places (clothing, floor) by a child aged at least 4~years, occurring at least once per month for at least three months, whether voluntary or involuntary.

\section*{Pathophysiology}
In the majority of cases (retentive encopresis), chronic constipation leads to fecal impaction with a large, hard stool mass that dilates the rectum. Liquid stool from the proximal colon then leaks around the impacted mass, resulting in soiling that the child may not sense due to chronic stretch-induced loss of rectal sensation.

\section*{Causes}
\begin{itemize}
  \item \textbf{Retentive (80--90\%):} Chronic functional constipation, painful defecation leading to withholding, stool toileting refusal
  \item \textbf{Non-retentive:} Emotional or behavioral disorders, oppositional defiant disorder, attention-deficit/hyperactivity disorder,
  \item \textbf{Anatomical:} Hirschsprung disease (rare, presents earlier), anal stenosis, anterior displaced anus
  \item \textbf{Neurologic:} Spinal cord abnormality, cerebral palsy, hypotonia
  \item \textbf{Psychosocial:} Toilet training conflicts, abuse, family stress, school bathroom avoidance
\end{itemize}

\section*{When to See a Doctor}
See your doctor if:
\begin{itemize}
  \item Failure to pass meconium in the first 48~hours of life (suggests Hirschsprung disease)
  \item Abdominal distension, vomiting, fever, or signs of obstruction
  \item Weight loss, failure to thrive, or neurologic deficits in lower extremities
\end{itemize}

\section*{Related Symptoms}
\begin{itemize}
  \item Chronic constipation, large-diameter stools, painful defecation
  \item Abdominal pain, bloating, decreased appetite
  \item Enuresis, urinary tract infections (stool mass compressing bladder)
\end{itemize}
\textbf{Important:} Encopresis is rarely a willful behavior; most children are embarrassed and distressed by it.

\section*{Self-Care / Home Management}
\begin{itemize}
  \item Establish a regular toileting schedule: sit on toilet for 5--10~minutes after meals
  \item Positive reinforcement for sitting and for successful bowel movements
  \item High-fiber diet (fruits, vegetables, whole grains) with adequate fluid intake
  \item Avoid punishment, shaming, or pressure; maintain a calm, supportive approach
\end{itemize}

\section*{How Your Doctor Will Evaluate You}
\begin{itemize}
  \item Abdominal palpation for fecal mass; rectal examination (under appropriate conditions) for impaction and tone
  \item Plain abdominal radiograph may show fecal loading in retentive cases
  \item Anorectal manometry and rectal biopsy if Hirschsprung disease is suspected
  \item Assess for comorbid emotional, behavioral, or learning difficulties
\end{itemize}

\section*{Treatment Options}
\begin{itemize}
  \item Disimpaction phase: polyethylene glycol (PEG 3350) 1--1.5~g/kg/day for 3--6~days
  \item Maintenance phase: PEG 3350 0.5--1~g/kg/day for at least 2--6~months to allow rectal tone to recover
  \item Behavioral modification with a reward system for toileting adherence
  \item Laxative weaning only after reliable, painless bowel movements have been established for weeks
\end{itemize}

\clearpage
